\documentclass[../main]{subfiles}
\begin{document}
\chapter{Aplicación de la Idea al Proyecto}
\actividad{Laboratorio de Ideas: Definición del Proyecto}

\begin{enumerate}
  \item \textbf{Organización:} Establecer los grupos de trabajo para
    el desarrollo del proyecto anual.
  \item \textbf{Investigación y Diagnóstico:}
    \begin{enumerate}
      \item Realizar relevamientos o encuestas para detectar fallas o
        necesidades en la comunidad o el taller.
      \item Priorizar los problemas encontrados segñun su
        viabilidad técnica.
    \end{enumerate}
  \item \textbf{Desarrollo de Soluciones:}
    \begin{enumerate}
      \item Proponer diferentes alternativas para los problemas detectados.
      \item Utilizar una matriz de ponderación para seleccionar la mejor opción.
      \item Definir la \textbf{Solución Base Optimizada} y proponer
        al menos 4 alternativas técnicas secundarias.
    \end{enumerate}
  \item \textbf{Formalización:} Redactar la idea de proyecto
    definitiva utilizando los elementos explicados en la teoría (Diagnóstico, Metas, Alcance, Recursos y Plan de Ejecución).
\end{enumerate}

\end{document}